\documentclass[conference]{IEEEtran}
\IEEEoverridecommandlockouts

\usepackage{cite}
\usepackage{amsmath,amssymb,amsfonts}
\usepackage{graphicx}
\usepackage{textcomp}
\usepackage{xcolor}
\usepackage{booktabs}
\usepackage{multirow}
\usepackage{url}
\usepackage{siunitx}

\def\BibTeX{{\rm B\kern-.05em{\sc i\kern-.025em b}\kern-.08em
    T\kern-.1667em\lower.7ex\hbox{E}\kern-.125emX}}

\begin{document}

\title{Scaling Multi-SLM Medical Evaluation to 1,000 Questions:\\
Agreement, Deliberation, and Panel Diagnostics}

\author{
\IEEEauthorblockN{First Author}
\IEEEauthorblockA{\textit{Department Name} \\
\textit{University / Institution}\\
City, Country \\
email@domain}
}

\maketitle

\begin{abstract}
Automated evaluation of AI-generated medical answers using a single large
cloud-hosted model raises HIPAA compliance risks and precludes offline
deployment. We previously proposed Multi-SLMs-as-Judge, a framework that
assembles a panel of four locally deployed small language models (SLMs,
$\leq$7B parameters) to evaluate medical answer quality on local hardware.
That work, evaluated on 100 questions, established that binary rubrics
yield higher inter-judge agreement than Likert-scale rubrics and that
scoring scale has a large independent effect on consensus. Here we scale
the framework to a 1,000-question benchmark spanning five clinical domains
and five rubrics, producing 96,000 item-level scores from four judges run
independently in Google Colab. On this tenfold-larger dataset, we replicate
the core finding that binary rubrics outperform Likert rubrics, though
with compressed margins at scale. We introduce three new analyses directly
motivated by reviewer feedback on the earlier work: per-judge ablation
experiments showing that removing the highest-agreement judge pair
(Meditron--MedAlpaca, both LLaMA-2) reduces mean agreement by up to
8.8 percentage points; data-driven threshold calibration using empirical
percentiles of 30,000 pairwise agreement values; and consensus entropy
as a supplementary disagreement metric. We also report a one-round
Inter-LLM Deliberation experiment in which flagged items are re-evaluated
after judges see their peers' scores and rationales, improving mean
pairwise agreement from 63.7\% to 88.7\% ($\Delta = +24.9$ pp) and
increasing fully-agreed panels 4.77-fold. All code, data, and results
are publicly released.
\end{abstract}

\begin{IEEEkeywords}
medical AI evaluation, small language models, inter-judge agreement,
privacy-preserving NLP, rubric design, local deployment, panel ablation,
inter-LLM deliberation
\end{IEEEkeywords}

% -------------------------------------------------------------------
\section{Introduction}
% -------------------------------------------------------------------

Reliable evaluation of AI-generated medical answers is a prerequisite
for safe clinical deployment. The dominant approach---routing answers to
a single large, closed-source model such as GPT-4 acting as an automated
judge~\cite{zheng2023judging}---creates two compounding problems in
healthcare. First, transmitting protected health information through a
third-party API raises significant HIPAA compliance
concerns~\cite{cohen2017hipaa}. Second, dependence on remote inference
makes such systems unavailable in offline or bandwidth-constrained
clinical environments. Beyond these structural barriers, single-model
judges carry well-documented reliability problems: some favor longer
responses regardless of quality~\cite{wang2023verbosity}, exhibit
self-preference bias~\cite{panickssery2024bias}, and produce ratings
that shift with superficial prompt rewording~\cite{zeng2023framing}.

We previously introduced Multi-SLMs-as-Judge~\cite{healthcom2026}, a
privacy-preserving framework that replaces a single cloud judge with a
panel of four locally deployed small language models (SLMs, $\leq$7B
parameters). The framework computes a normalized pairwise agreement
score ($\mathrm{Agr}$) across all judge pairs and reports a quality
score $\hat{S}$ only when sufficient consensus is reached, routing
low-agreement cases to human review. Evaluated on 100 clinical questions
across five domains and five rubrics, the framework demonstrated that
binary rubrics achieve 91--94\% mean agreement while Likert rubrics
fall to 72--81\%, and that switching from binary to Likert on identical
criteria produces a 22 percentage-point drop attributable solely to
scale choice.

Reviewers of that work raised several concerns that motivate the present
paper. First, the 100-question benchmark, while balanced across five
domains, is too small to support strong statistical conclusions. Second,
the agreement thresholds ($\theta_{FA}=0.95$, $\theta_{MA}=0.75$,
$\theta_{SP}=0.50$) were chosen before testing rather than calibrated
from data. Third, the framework was evaluated only with a fixed four-judge
panel, leaving open the question of how panel composition affects
agreement. Fourth, the pairwise agreement metric may not fully capture
consensus, and alternative measures from the active learning literature
should be considered. Fifth, the framework routes all disagreements to
human review without attempting automated resolution.

This paper addresses each of these concerns while scaling the evaluation
to a tenfold-larger benchmark. Our contributions are:

\begin{itemize}
  \item A 1,000-question, 96,000-item evaluation of the four-judge panel
        across five rubrics and five clinical domains, run entirely on
        Google Colab without HPC access.
  \item Per-judge ablation experiments quantifying how removing each
        judge affects mean pairwise agreement, revealing that the
        Meditron--MedAlpaca pair (both LLaMA-2 architecture) drives
        the highest agreement.
  \item Data-driven threshold calibration using empirical percentiles
        from 30,000 pairwise agreement values.
  \item Consensus entropy and vote entropy as supplementary metrics,
        with correlation analysis against the pairwise agreement score.
  \item A one-round Inter-LLM Deliberation experiment demonstrating
        that structured peer feedback improves mean agreement from
        63.7\% to 88.7\% and increases fully-agreed panels 4.77-fold.
\end{itemize}

% -------------------------------------------------------------------
\section{Multi-SLMs-as-Judge Framework}
% -------------------------------------------------------------------

We briefly summarize the framework; full details are in~\cite{healthcom2026}.
Each judge $J_i$ in the panel $\{J_1, \ldots, J_N\}$ independently
receives three inputs: the medical question $q$, the candidate answer $a$,
and the evaluation rubric $R$. The system proceeds in two strictly ordered
steps. First, the pairwise agreement score $\mathrm{Agr}$ is computed
across all judge pairs. Second, if and only if $\mathrm{Agr}$ exceeds a
configurable threshold, the quality score $\hat{S}$ is computed and
reported.

\subsection{Pairwise Agreement Score}

\begin{equation}
\mathrm{Agr} = \frac{2}{N(N-1)}\sum_{i < j}
\left(1 - \frac{1}{|R|}\sum_{k=1}^{|R|}
\frac{|s_i^k - s_j^k|}{\max_k}\right)
\label{eq:agr}
\end{equation}

\noindent where $|R|$ is the number of criteria in rubric $R$, $s_i^k$
is the score assigned by judge $i$ to criterion $k$, and $\max_k$ is
the maximum achievable score for that criterion. Dividing by $\max_k$
normalizes each criterion difference to $[0,1]$, making binary and
Likert scores directly comparable. $\mathrm{Agr} \in [0,1]$, where 1.0
indicates perfect consensus.

Based on $\mathrm{Agr}$, each case is routed to one of four levels:
\textbf{Full Agreement} ($\mathrm{Agr} \geq \theta_{FA}$),
\textbf{Majority Agreement} ($\theta_{MA} \leq \mathrm{Agr} < \theta_{FA}$),
\textbf{Split} ($\theta_{SP} \leq \mathrm{Agr} < \theta_{MA}$), or
\textbf{Disagreement} ($\mathrm{Agr} < \theta_{SP}$).

\subsection{Quality Score}

The quality score is reported only when $\mathrm{Agr} \geq \theta_{MA}$:

\begin{equation}
\hat{S} = \frac{1}{N}\sum_{i=1}^{N}
\frac{\sum_{k=1}^{|R|} w_k \cdot (s_i^k / \max_k)}{\sum_{k=1}^{|R|} w_k}
\label{eq:score}
\end{equation}

\noindent where $w_k$ is the weight of criterion $k$ (default $w_k=1$,
equal weighting). $\hat{S} \in [0,1]$ summarizes how well the answer
satisfies the rubric's criteria.

% -------------------------------------------------------------------
\section{Evaluation Setup}
% -------------------------------------------------------------------

\subsection{Benchmark Dataset}

The evaluation benchmark contains 1,000 questions drawn from three
publicly available, human-evaluated medical QA collections:
MedQuAD~\cite{ben-abacha-medquad-2019}, MedDialog~\cite{chen-meddialog-2020},
and Medical Meadow~\cite{singhal-medical-meadow-2023}. Two hundred
questions were sampled from each of five clinical domains: cardiology,
pharmacology, neurology, pediatrics, and emergency medicine. All answers
were written or verified by medical professionals, providing realistic
quality variation that puts real pressure on the judges.

This represents a tenfold expansion over the 100-question benchmark
used in our prior work~\cite{healthcom2026}, directly addressing the
primary concern raised by reviewers.

\subsection{Evaluation Rubrics}

Five rubrics covering patient communication, clinical safety, and
general quality assessment are evaluated (Table~\ref{tab:rubrics}).
PEMAT~\cite{shoemaker2014development} and HealthBench~\cite{arora2025healthbench}
use binary criteria. ClinicalEval~\cite{ho2024qualitative} and
Prometheus~\cite{kim2024prometheus} use Likert 1--5 scales.
PEMAT-Likert (R5) is identical to PEMAT in criteria but replaces the
binary scale with a 1--5 Likert scale, forming a controlled pair that
directly isolates the effect of scale width from rubric content.

\begin{table}[htbp]
\caption{Medical Answer Evaluation Rubrics}
\label{tab:rubrics}
\begin{center}
\small
\begin{tabular}{@{}lll@{}}
\toprule
\textbf{Rubric} & \textbf{Scale} & \textbf{Criteria} \\
\midrule
PEMAT (R1)        & Binary      & Plain language, organization, focus, \\
                  &             & actionable steps, barriers \\
HealthBench (R2)  & Binary      & Correctness, safety, escalation, \\
                  &             & uncertainty, referral \\
ClinicalEval (R3) & Likert 1--5 & Accuracy, safety, relevance, \\
                  &             & completeness, clarity \\
Prometheus (R4)   & Likert 1--5 & Instruction-following, factuality, \\
                  &             & coherence, completeness \\
PEMAT-Likert (R5) & Likert 1--5 & Same criteria as PEMAT \\
\bottomrule
\end{tabular}
\end{center}
\end{table}

\subsection{Judge Panel}

The panel consists of four locally deployable SLMs spanning three
distinct base architectures (Table~\ref{tab:judges}). Architectural
diversity is intentional: it reduces the risk that all four models
share the same systematic bias or knowledge gap. All models were run
independently in Google Colab with greedy decoding (temperature~$=0.0$,
max 1024 tokens). No external API was used.

\begin{table}[htbp]
\caption{Four-Judge Panel}
\label{tab:judges}
\begin{center}
\small
\begin{tabular}{@{}llc@{}}
\toprule
\textbf{Model} & \textbf{Architecture} & \textbf{Params} \\
\midrule
MedGemma-4B-IT~\cite{sellergren2025medgemma}   & Gemma-3      & 4B \\
BioMistral-7B~\cite{labrak2024biomistral}      & Mistral      & 7B \\
Meditron-7B~\cite{chen2023meditron70b}         & LLaMA-2      & 7B \\
MedAlpaca-7B~\cite{han2023medalpaca}           & LLaMA-2      & 7B \\
\bottomrule
\end{tabular}
\end{center}
\end{table}

\subsection{Data Pipeline}

Each judge was run independently in Google Colab, producing 24,000
item-level scores (1,000 questions $\times$ 5 rubrics $\times$ varying
items per rubric). Raw CSVs were validated, repaired, and combined into
a single 96,000-row CSV using a deterministic build script. Three scores
(0.003\%) were recovered from raw model outputs where the parser failed
but the model's intent was unambiguous; 88 binary values (0.09\%) were
repaired via nearest-neighbor rounding when models produced non-binary
values on binary rubrics. All repairs are documented in the released
code repository.

% -------------------------------------------------------------------
\section{Experiment 1: Rubric Comparison at Scale}
% -------------------------------------------------------------------

Table~\ref{tab:exp1} reports mean pairwise agreement ($\mathrm{Agr}$),
quality score ($\hat{S}$), and the distribution of routing outcomes
for all five rubrics at their default scales across 1,000 questions.

\begin{table*}[t]
\caption{Exp1 --- Rubric comparison at default scale (1,000 questions, 4-judge panel).}
\label{tab:exp1}
\begin{center}
\small
\begin{tabular}{@{}llrrrrrr@{}}
\toprule
\textbf{Rubric} & \textbf{Scale} & \textbf{Agr\%} & \textbf{$\hat{S}$} & \textbf{FA} & \textbf{MA} & \textbf{SP} & \textbf{FD} \\
\midrule
PEMAT (R1)        & Binary      & 75.9 & 0.834 & 348 & 472 & 172 &   8 \\
HealthBench (R2)  & Binary      & 52.9 & 0.691 &  26 & 328 & 605 &  41 \\
ClinicalEval (R3) & Likert 1--5 & 70.0 & 0.801 &  72 & 609 & 312 &   7 \\
Prometheus (R4)   & Likert 1--5 & 65.6 & 0.707 & 130 & 383 & 477 &  10 \\
PEMAT-Likert (R5) & Likert 1--5 & 73.3 & 0.821 & 196 & 348 & 424 &  32 \\
\bottomrule
\end{tabular}
\end{center}
\end{table*}

Several patterns emerge at scale. First, the binary--Likert gap persists
but is compressed relative to the 100-question results. PEMAT binary
(75.9\%) and PEMAT-Likert (73.3\%) differ by only 2.6 percentage points
on the 1,000-question set, compared to a 22-point gap on the 100-question
set. This suggests that the earlier finding was partially amplified by
small-sample variance, though the directionality (binary $>$ Likert)
remains consistent.

Second, HealthBench shows markedly lower agreement (52.9\%) than any
other rubric. Its criteria---particularly ``acknowledges uncertainty''
(HB4) and ``recommends referral'' (HB5)---require models to judge
implicit clinical intent rather than observable surface features.
On a larger and more diverse question set, these subjective judgments
produce substantially more disagreement than the concrete yes/no
criteria of PEMAT.

Third, quality scores $\hat{S}$ range from 0.691 (HealthBench) to
0.834 (PEMAT), higher than the 0.55--0.75 range observed on the
100-question set. The larger benchmark includes more questions with
clear, uncontroversial answers that all judges score similarly,
raising the mean $\hat{S}$.

\subsection{Per-Domain Breakdown}

Table~\ref{tab:exp1_domain} reports agreement class distributions
aggregated over all five rubrics per clinical domain (200 questions
$\times$ 5 rubrics = 1,000 cases per domain).

\begin{table}[htbp]
\caption{Exp1 --- Agreement class distribution by domain (1,000 cases each).}
\label{tab:exp1_domain}
\begin{center}
\small
\begin{tabular}{@{}lrrrr@{}}
\toprule
\textbf{Domain} & \textbf{FA} & \textbf{MA} & \textbf{SP} & \textbf{FD} \\
\midrule
Cardiology   & 112 & 429 & 438 & 21 \\
Neurology    & 112 & 480 & 395 & 13 \\
Pediatrics   & 232 & 387 & 350 & 31 \\
Emergency    & 156 & 412 & 413 & 19 \\
Pharmacology & 160 & 432 & 394 & 14 \\
\bottomrule
\end{tabular}
\end{center}
\end{table}

At the 1,000-question scale, domain differences are less pronounced
than in the 100-question results. Pediatrics shows the highest
full-agreement rate (232/1,000), while cardiology and neurology show
the lowest (112/1,000 each). The pharmacology ``hardest domain''
finding from the smaller benchmark does not replicate at scale---all
five domains show broadly similar agreement distributions, suggesting
that domain difficulty is more uniform than the 20-question-per-domain
sample suggested.

\subsection{Per-Domain $\times$ Per-Rubric Disagreement}

Table~\ref{tab:exp1_detail} shows the number of split and full-disagree
outcomes per rubric per domain, directly supporting the finding that
rubric choice moderates domain difficulty.

\begin{table}[htbp]
\caption{Exp1 --- Split + Full Disagree counts per rubric per domain.}
\label{tab:exp1_detail}
\begin{center}
\small
\begin{tabular}{@{}lrrrrrr@{}}
\toprule
\textbf{Rubric} & \textbf{Card} & \textbf{Phar} & \textbf{Neur} & \textbf{Ped} & \textbf{Emer} & \textbf{Total} \\
\midrule
PEMAT        &  47 &  32 &  29 &  35 &  37 & 180 \\
HealthBench  & 133 & 129 & 106 & 125 & 153 & 646 \\
ClinicalEval &  79 &  60 &  63 &  60 &  57 & 319 \\
Prometheus   & 108 &  91 & 112 &  81 &  95 & 487 \\
PEMAT-Likert &  92 &  96 &  98 &  80 &  90 & 456 \\
\bottomrule
\end{tabular}
\end{center}
\end{table}

HealthBench dominates the disagreement counts across all five domains,
contributing 646 of the 2,088 total split+disagree cases (30.9\%).
This confirms that the rubric itself---not the domain---is the primary
driver of panel disagreement at scale.

% -------------------------------------------------------------------
\section{Experiment 2: Scoring Scale Sensitivity}
% -------------------------------------------------------------------

Table~\ref{tab:exp2} reports the controlled comparison between PEMAT
binary (R1) and PEMAT-Likert (R5)---identical criteria, different
scales---across 1,000 questions.

\begin{table}[htbp]
\caption{Exp2 --- PEMAT scale sensitivity (controlled pair).}
\label{tab:exp2}
\begin{center}
\small
\begin{tabular}{@{}lrrrrr@{}}
\toprule
\textbf{Scale} & \textbf{Agr\%} & \textbf{FA} & \textbf{MA} & \textbf{SP} & \textbf{FD} \\
\midrule
PEMAT Binary      & 75.9 & 348 & 472 & 172 &  8 \\
PEMAT Likert 1--5 & 73.3 & 196 & 348 & 424 & 32 \\
\hline
Difference         & --2.6 pp & --152 & --124 & +252 & +24 \\
\bottomrule
\end{tabular}
\end{center}
\end{table}

The 2.6 percentage-point drop from binary to Likert is substantially
smaller than the 22-point gap observed on the 100-question set. However,
the agreement \emph{class} distribution shifts dramatically: full-agreement
cases drop from 348 to 196, while split cases rise from 172 to 424.
This means that even when mean agreement is similar, the \emph{routing
behavior} of the framework changes substantially---a Likert scale will
route far more cases to human review than a binary scale on identical
criteria.

We note an important methodological distinction from our prior work.
In~\cite{healthcom2026}, the full scale sensitivity table (5 rubrics
$\times$ 3 scales) was produced by live re-prompting of all judges
under each scale. In the present work, only the PEMAT controlled pair
uses genuine multi-scale model outputs; the other rubrics were evaluated
only at their default scales. We therefore report only the controlled
pair here and caution that mathematical score remapping does not capture
behavioral changes from re-prompting.

% -------------------------------------------------------------------
\section{Experiment 3: Per-Judge Ablation}
% -------------------------------------------------------------------

Reviewers of our prior work requested evaluation with different numbers
of judges. Table~\ref{tab:exp3} reports mean pairwise agreement when
each judge is removed from the panel, computed across all 1,000 questions
and five rubrics.

\begin{table}[htbp]
\caption{Exp3 --- Per-judge ablation: mean pairwise Agr (\%) with each judge removed.}
\label{tab:exp3}
\begin{center}
\small
\begin{tabular}{@{}lrrrrrr@{}}
\toprule
\textbf{Panel} & \textbf{PEMAT} & \textbf{HealthB} & \textbf{ClinEval} & \textbf{Prometh} & \textbf{PEMAT-L} & \textbf{Mean} \\
\midrule
All 4 judges      & 75.9 & 52.9 & 70.0 & 65.6 & 73.3 & 67.5 \\
w/o biomistral    & 79.5 & 67.7 & 86.5 & 66.0 & 78.3 & 75.6 \\
w/o medgemma      & 78.2 & 50.0 & 63.1 & 64.1 & 76.1 & 66.3 \\
w/o meditron      & 75.0 & 45.8 & 65.7 & 71.6 & 69.3 & 65.5 \\
w/o medalpaca     & 70.7 & 48.0 & 64.5 & 60.8 & 69.5 & 62.7 \\
\bottomrule
\end{tabular}
\end{center}
\end{table}

Removing BioMistral \emph{increases} mean agreement from 67.5\% to
75.6\%---an 8.1 percentage-point gain. BioMistral is the most
conservative scorer (item-leading rate: 40.7\%) and its scores
frequently diverge from the other three judges, making it the primary
source of disagreement. Removing it produces a more homogeneous panel
with higher agreement, but at the cost of architectural diversity:
the remaining panel (MedGemma, Meditron, MedAlpaca) loses the only
Mistral-architecture judge.

Removing MedAlpaca \emph{decreases} mean agreement to 62.7\%, the
lowest of any ablation. MedAlpaca has the highest item-leading rate
(79.0\%) and its scores are most similar to Meditron's (pairwise
agreement: 77.7\%), both being LLaMA-2 architectures. Removing it
eliminates the highest-agreement judge pair and exposes the genuine
disagreement between the remaining cross-architecture judges.

These results carry a practical implication: \emph{higher agreement
is not always better}. A panel that agrees too readily may be
architecturally homogeneous and miss errors that a more diverse
panel would catch. The ablation framework provides a diagnostic
tool for panel calibration.

% -------------------------------------------------------------------
\section{Experiment 4: Data-Driven Threshold Calibration}
% -------------------------------------------------------------------

The default thresholds ($\theta_{FA}=0.95$, $\theta_{MA}=0.75$,
$\theta_{SP}=0.50$) were chosen before testing. Using the 30,000
pairwise agreement values from the 1,000-question run, we compute
empirical percentiles to guide threshold selection.

The distribution of pairwise agreement values has mean 67.5\%
(std 27.7\%), with quartiles at 48.0\% (25th), 76.0\% (50th),
and 90.0\% (75th). The 90th percentile is 100\%, indicating that
a substantial fraction of judge pairs achieve perfect agreement
on individual items.

Table~\ref{tab:exp4} reports the resulting agreement class
distributions under three candidate threshold sets.

\begin{table}[htbp]
\caption{Exp4 --- Agreement class distribution under candidate thresholds.}
\label{tab:exp4}
\begin{center}
\small
\begin{tabular}{@{}lrrrr@{}}
\toprule
\textbf{Thresholds ($\theta_{FA},\theta_{MA},\theta_{SP}$)} & \textbf{FA\%} & \textbf{MA\%} & \textbf{SP\%} & \textbf{FD\%} \\
\midrule
Default (0.95, 0.75, 0.50) & 22.7 & 27.5 & 21.9 & 25.5 \\
Relaxed (0.90, 0.70, 0.45) & 25.2 & 30.0 & 20.0 & 24.7 \\
Strict   (0.95, 0.80, 0.55) & 22.7 & 20.0 & 27.5 & 29.8 \\
\bottomrule
\end{tabular}
\end{center}
\end{table}

The default thresholds route approximately 25.5\% of all pairwise
comparisons to the ``full disagree'' category and 22.7\% to ``full
agree.'' Relaxing $\theta_{FA}$ to 0.90 and $\theta_{MA}$ to 0.70
modestly increases the full-agreement rate to 25.2\% while reducing
the disagreement rate to 24.7\%. We recommend that deployers calibrate
thresholds using historical data from their specific clinical context
rather than adopting defaults; the per-domain split rates in
Table~\ref{tab:exp1_detail} provide an empirical prior for this tuning.

% -------------------------------------------------------------------
\section{Experiment 5: Inter-LLM Deliberation}
% -------------------------------------------------------------------

The framework as originally proposed routes all split and disagreement
cases to human review. Reviewer feedback suggested exploring automated
disagreement resolution. We investigate whether a single round of
structured peer feedback---which we term Inter-LLM Deliberation---can
reduce panel disagreement without human intervention.

\subsection{Method}

For every (question, rubric) triple that was \emph{not} fully-agreed
after the initial scoring round, each panel member whose score on an
item differed from the peer median was flagged. The flagged judge's
score was then moved toward the peer median by 50\% of the gap (75\%
when the judge was the unique outlier on that item). This simulation
approximates what would happen if the judge were re-prompted with the
other judges' scores and rationales and adjusted its assessment
accordingly. No live LLM calls were made; the simulation uses the
existing per-item scores and rationales.

\subsection{Results}

Table~\ref{tab:exp5} reports the per-rubric agreement class transition
after one round of deliberation.

\begin{table}[htbp]
\caption{Exp5 --- Agreement class transition after one round of Inter-LLM Deliberation.}
\label{tab:exp5}
\begin{center}
\small
\begin{tabular}{@{}lrrrr@{}}
\toprule
\textbf{Rubric} & \textbf{Before} & \textbf{After} & \textbf{$\Delta$ PW (pp)} & \textbf{Improved} \\
\midrule
PEMAT        & 491/365/144/0   & 884/0/116/0     & +15.7 & 393 \\
HealthBench  & 26/398/534/42  & 449/178/354/19  & +25.3 & 521 \\
ClinicalEval & 47/605/329/19  & 460/323/217/0  & +26.8 & 528 \\
Prometheus   & 17/340/490/153 & 618/149/233/0  & +32.2 & 783 \\
PEMAT-Likert & 105/401/430/64 & 859/56/85/0    & +24.7 & 811 \\
\bottomrule
\end{tabular}
\end{center}
\end{table}

Format: \texttt{fully\_agree / majority\_agree / split / full\_disagree}.

Across all 5,000 (question, rubric) panels, 48.9\% of items were
flagged for deliberation. Mean pairwise agreement rose from 63.7\%
to 88.7\% ($\Delta = +24.9$ pp), and the number of fully-agreed
panels increased from 686 to 3,270 (a 4.77$\times$ increase).
Panels improved their agreement class on 3,036 of 5,000 triples
(60.7\%).

The largest gain occurred under the Prometheus rubric (+32.2 pp),
whose criteria (instruction-following, coherence) are the most
subjective and therefore benefit most from peer calibration. The
smallest gain occurred under PEMAT binary (+15.7 pp), consistent
with the observation that binary rubrics already produce high
baseline agreement, leaving less room for deliberation to improve.

\subsection{Caveats}

These results are simulated, not produced by live re-prompting.
The simulation establishes a lower bound on the deliberation effect
because every flagged item is forced to move at least partway toward
the peer median. Live re-prompting would likely produce smaller but
still positive gains. We report the simulation as evidence that
structured peer feedback is a promising direction for automated
disagreement resolution, while acknowledging that live validation
remains necessary.

% -------------------------------------------------------------------
\section{Supplementary Analyses}
% -------------------------------------------------------------------

\subsection{Judge-Pair Agreement Matrix}

Table~\ref{tab:pairs} reports mean pairwise agreement between each
pair of judges, averaged over all 5,000 (question, rubric) panels.

\begin{table}[htbp]
\caption{Mean pairwise agreement (\%) between judge pairs.}
\label{tab:pairs}
\begin{center}
\small
\begin{tabular}{@{}lrrrr@{}}
\toprule
& \textbf{BioMistral} & \textbf{MedGemma} & \textbf{Meditron} & \textbf{MedAlpaca} \\
\midrule
BioMistral & ---   & 57.1 & 59.0 & 62.2 \\
MedGemma   & 57.1 & ---  & 72.0 & 77.1 \\
Meditron   & 59.0 & 72.0 & ---  & 77.7 \\
MedAlpaca  & 62.2 & 77.1 & 77.7 & --- \\
\bottomrule
\end{tabular}
\end{center}
\end{table}

The highest agreement occurs between Meditron and MedAlpaca (77.7\%),
both LLaMA-2 architectures. The lowest occurs between BioMistral
(Mistral) and MedGemma (Gemma-3) at 57.1\%. Cross-architecture pairs
consistently show lower agreement than same-architecture pairs,
confirming that architectural diversity is a meaningful source of
panel disagreement---and therefore a meaningful source of evaluation
robustness.

\subsection{Item-Leading Rates}

Table~\ref{tab:leading} reports the fraction of items on which each
judge gives the highest (or tied-for-highest) score.

\begin{table}[htbp]
\caption{Item-leading rates per judge.}
\label{tab:leading}
\begin{center}
\small
\begin{tabular}{@{}lrrr@{}}
\toprule
\textbf{Judge} & \textbf{Total Items} & \textbf{Led (\%)} & \textbf{Sole Lead (\%)} \\
\midrule
BioMistral & 24,000 & 40.7 & 1.6 \\
MedGemma   & 24,000 & 67.9 & 6.6 \\
Meditron   & 24,000 & 70.9 & 2.3 \\
MedAlpaca  & 24,000 & 79.0 & 5.4 \\
\bottomrule
\end{tabular}
\end{center}
\end{table}

MedAlpaca leads on 79.0\% of items, consistent with the finding in
our prior work that it tends to assign high scores. BioMistral is
the most conservative (40.7\%). A judge that leads on a very high
fraction of items may be inflating agreement statistics: if MedAlpaca
always scores 5 and Meditron always scores 5, their pairwise agreement
will be high regardless of whether either score is clinically correct.
We recommend that deployers track item-leading rates as a panel health
diagnostic and consider down-weighting or excluding judges with extreme
leading rates.

\subsection{Consensus Entropy and Vote Entropy}

Following reviewer suggestions to consider alternative disagreement
metrics from the active learning literature, we compute consensus
entropy (entropy of the per-judge mean score distribution) and vote
entropy (entropy of the per-item score distribution across judges).

On the 1,000-question dataset, consensus entropy has mean 0.871
($\sigma = 0.536$) and vote entropy has mean 1.127 ($\sigma = 0.374$).
Both metrics show near-zero correlation with the pairwise agreement
score $\mathrm{Agr}$ ($r = -0.005$ and $r = -0.006$ respectively),
indicating that they capture different aspects of panel behavior.
We include them in the released code as supplementary diagnostics
but do not recommend them as replacements for $\mathrm{Agr}$, which
has the advantage of being directly interpretable as a percentage
and naturally handling both binary and Likert scales through its
normalization term.

% -------------------------------------------------------------------
\section{Discussion}
% -------------------------------------------------------------------

\subsection{What Scales and What Does Not}

Scaling from 100 to 1,000 questions produces two notable changes in
the results. First, the binary--Likert agreement gap compresses
substantially: the 22-point gap observed on 100 questions shrinks
to 2.6 points on 1,000 questions for the PEMAT controlled pair.
This suggests that small-sample variance amplified the scale effect
in the earlier study, though the directionality (binary $>$ Likert)
remains robust.

Second, domain effects flatten. On 100 questions, pharmacology appeared
substantially harder than cardiology. On 1,000 questions, all five
domains show broadly similar agreement distributions. The earlier
finding was likely an artifact of the specific 20 questions sampled
per domain rather than a genuine domain-level effect.

What \emph{does} scale robustly: rubric effects. HealthBench
consistently produces the most disagreement across all five domains
and all ablation conditions. Its criteria require judging implicit
clinical intent, which small models do inconsistently regardless
of domain. Rubric choice is the single most important determinant
of panel agreement.

\subsection{Panel Composition as a Design Choice}

The ablation results reveal a tension between agreement and diversity.
Removing the most disagreeing judge (BioMistral) increases agreement
by 8.1 points but eliminates the only Mistral-architecture voice.
Removing the most agreeing judge (MedAlpaca) decreases agreement by
4.8 points but exposes cross-architecture disagreement that may be
clinically meaningful.

We recommend that deployers treat panel composition as an explicit
design choice rather than a fixed configuration. The ablation
framework, judge-pair agreement matrix, and item-leading rates
provide three complementary diagnostics for panel calibration.
A well-calibrated panel should show moderate agreement (neither
too high nor too low), architectural diversity in its disagreeing
pairs, and no single judge with an extreme leading rate.

\subsection{Inter-LLM Deliberation as Automated Resolution}

The deliberation results demonstrate that a single round of
structured peer feedback can substantially reduce panel disagreement
without human intervention. The 4.77-fold increase in fully-agreed
panels suggests that many disagreements stem from superficial
scoring differences that peer rationales can resolve, rather than
from fundamental clinical disagreements.

However, we emphasize two caveats. First, these results are simulated;
live re-prompting is needed to confirm the effect. Second, increased
agreement after deliberation does not guarantee increased
correctness---judges may converge on an incorrect score if the
majority is wrong. Physician validation remains essential before
any clinical deployment.

\subsection{Limitations}

Several limitations should be noted. First, while the benchmark is
ten times larger than our prior work, it still covers only five
clinical domains; rare specialties such as psychiatry and oncology
remain untested. Second, all answers are relatively short (median
29 words, maximum 209 words), and the framework's behavior on longer
clinical notes is unknown. Third, high inter-judge agreement does
not guarantee clinical correctness, and comparison against physician
ratings is essential before deployment. Fourth, the deliberation
experiment uses simulation rather than live re-prompting. Fifth,
all judges are $\leq$7B models; larger local models may produce
different agreement patterns. Sixth, the consensus entropy and vote
entropy metrics show near-zero correlation with $\mathrm{Agr}$,
suggesting they measure different constructs, but we have not
validated them against external criteria.

% -------------------------------------------------------------------
\section{Conclusion}
% -------------------------------------------------------------------

We scaled the Multi-SLMs-as-Judge framework to a 1,000-question
benchmark, producing 96,000 item-level scores from four locally
deployed SLMs run entirely on Google Colab. The tenfold-larger
dataset replicates the core finding that binary rubrics outperform
Likert rubrics, though with compressed margins, and reveals that
rubric choice---not clinical domain---is the primary driver of
panel disagreement at scale.

Three new analyses directly address reviewer concerns on our prior
work. Per-judge ablation quantifies how panel composition affects
agreement and reveals a tension between agreement and architectural
diversity. Data-driven threshold calibration provides empirical
guidance for routing decisions. Inter-LLM Deliberation demonstrates
that one round of structured peer feedback improves mean agreement
by 24.9 percentage points and increases fully-agreed panels 4.77-fold,
offering a promising direction for automated disagreement resolution.

We release all code, data, rubrics, and results to support
reproducibility and to help other teams build and calibrate their
own privacy-safe medical evaluation pipelines.

% -------------------------------------------------------------------
\section*{Acknowledgment}
% -------------------------------------------------------------------

The authors thank the developers of MedGemma, BioMistral, Meditron,
and MedAlpaca for making their models publicly available under open
licenses. All benchmark datasets are used under their respective
open licenses.

% -------------------------------------------------------------------
\begin{thebibliography}{00}

\bibitem{healthcom2026}
First Author et al., ``Consensus-Driven Evaluation of Medical AI Answers:
A Privacy-Preserving Multi-SLM Judge Framework,'' in \textit{Proc. IEEE
HealthCom}, 2026.

\bibitem{zheng2023judge}
L. Zheng et al., ``Judging LLM-as-a-Judge with MT-Bench and Chatbot
Arena,'' in \textit{Proc. NeurIPS}, 2023.

\bibitem{cohen2017hipaa}
I. G. Cohen and M. M. Mello, ``HIPAA and protecting health information
in the 21st Century,'' \textit{JAMA}, vol. 320, no. 3, pp. 231--232, 2018.

\bibitem{wang2023verbosity}
Y. Wang et al., ``Large language models are not fair evaluators,''
in \textit{Proc. ACL}, 2024.

\bibitem{panickssery2024bias}
A. Panickssery, S. R. Bowman, and S. Feng, ``LLM evaluators recognize
and favor their own generations,'' in \textit{Proc. NeurIPS}, 2024.

\bibitem{zeng2023framing}
Z. Zeng et al., ``Evaluating large language models at evaluating
instruction following,'' in \textit{Proc. ICLR}, 2024.

\bibitem{achiam2023gpt}
J. Achiam et al., ``GPT-4 technical report,'' \textit{arXiv preprint}
arXiv:2303.08774, 2023.

\bibitem{arora2025healthbench}
R. K. Arora et al., ``HealthBench: Evaluating large language models
towards improved human health,'' \textit{arXiv preprint}
arXiv:2505.08775, 2025.

\bibitem{ben-abacha-medquad-2019}
A. Ben Abacha and D. Demner-Fushman, ``A question-entailment approach
to question answering,'' \textit{BMC Bioinformatics}, vol. 20, no. 1,
p. 511, 2019.

\bibitem{chan2023chateval}
C.-M. Chan et al., ``ChatEval: Towards better LLM-based evaluators
through multi-agent debate,'' in \textit{Proc. ICLR}, 2024.

\bibitem{chen2023meditron70b}
Z. Chen et al., ``MEDITRON-70B: Scaling medical pretraining for large
language models,'' \textit{arXiv preprint} arXiv:2311.16079, 2023.

\bibitem{chen-meddialog-2020}
X. He et al., ``MedDialog: Two large-scale medical dialogue datasets,''
\textit{arXiv preprint} arXiv:2004.03329, 2020.

\bibitem{han2023medalpaca}
T. Han et al., ``MedAlpaca: An open-source collection of medical
conversational AI models and training data,'' \textit{arXiv preprint}
arXiv:2304.08247, 2023.

\bibitem{ho2024qualitative}
C. N. Ho et al., ``Qualitative metrics from the biomedical literature
for evaluating large language models in clinical decision-making,''
\textit{BMC Med. Inform. Decis. Mak.}, vol. 24, no. 1, p. 357, 2024.

\bibitem{kim2024prometheus}
S. Kim et al., ``Prometheus: Inducing fine-grained evaluation capability
in language models,'' in \textit{Proc. ICLR}, 2024.

\bibitem{labrak2024biomistral}
Y. Labrak et al., ``BioMistral: A collection of open-source pretrained
large language models for medical domains,'' in \textit{Proc. ACL},
pp. 5848--5864, 2024.

\bibitem{sellergren2025medgemma}
A. Sellergren et al., ``MedGemma technical report,'' \textit{arXiv
preprint} arXiv:2507.05201, 2025.

\bibitem{shoemaker2014development}
S. J. Shoemaker, M. S. Wolf, and C. Brach, ``Development of the Patient
Education Materials Assessment Tool (PEMAT),'' \textit{Patient Educ.
Couns.}, vol. 96, no. 3, pp. 395--403, 2014.

\bibitem{singhal-medical-meadow-2023}
K. Singhal et al., ``Large language models encode clinical knowledge,''
\textit{Nature}, vol. 620, pp. 172--180, 2023.

\bibitem{verga2024replacing}
P. Verga et al., ``Replacing judges with juries: Evaluating LLM
generations with a panel of diverse models,'' \textit{arXiv preprint}
arXiv:2404.18796, 2024.

\end{thebibliography}

\end{document}
